% ============================ 3. CASE STUDIES ==================================
\newpage
\phase{Section three}{How ten companies grew, or did not}{Four Indian developers, two Indian technology start-ups, three global technology companies and one cautionary tale.}

\subsection{The Indian developers: power first, molecules second}

{\tablefont
\begin{xltabular}{\textwidth}{@{}>{\RaggedRight\arraybackslash}p{18mm} L L L L@{}}
\caption{Indian developers. All figures \rep\ unless tagged otherwise.}\\
\toprule
\thead{Company} & \thead{Founding and capital} & \thead{First hydrogen step and inflection} & \thead{Where it stands (2025–26)} & \thead{Mistake or lesson}\\ \midrule
\endfirsthead
\toprule
\thead{Company} & \thead{Founding and capital} & \thead{First hydrogen step and inflection} & \thead{Where it stands (2025–26)} & \thead{Mistake or lesson}\\ \midrule
\endhead
\textbf{ACME} & Founded by Manoj Upadhyay \cite{acmeihi2024}; ACME Solar Holdings incorporated 2015; IPO of ₹2,900 crore in Nov 2024 (₹2,395 crore fresh) \cite{acmeipo2024} & Built what may be the world's first green-ammonia plant at Bikaner in 2021 \cite{acmeihi2024}. Inflection: won 370,000 t/yr, 51\% of SECI's Tranche I \cite{seci2025} & FY25 revenue ₹1,575 crore, profit ₹251 crore, 2,705 MW operating, 6,970 MW portfolio \cite{acmefy25}. Gopalpur: 0.4 Mt/yr to Japan's IHI, first output likely 2027 \cite{acmeihi2024} \prj & Built a cash-generating power business first, then used it as the platform. Risk: half the national auction volume rides on one bidder's cost assumptions.\\
\textbf{Greenko / AM Green} & Founder Anil Chalamalasetty \cite{amgreenfid2024}. GIC (majority) and ADIA added \usd{}329 m in 2019, part of \usd{}824 m committed \cite{greenkogic2019}. In 2025 AM Green bought ORIX's 17.5\% of Greenko, taking its group to 25\% \cite{greenkoorix2025} & FID in 2024 on a 1 Mt/yr green-ammonia plant at Kakinada, firmed by 950 MW of pumped storage, with John Cockerill alkaline electrolysers \cite{amgreenfid2024}. Inflection: Uniper signs for up to 500,000 t/yr (Jan 2026) \cite{amgreenuniper2026} & Over 10 GW of renewables and pumped storage \cite{greenkoorix2025}. First shipment to Uniper ``as early as 2028'' \cite{amgreenuniper2026} \prj & Pumped storage gives it the steadiest power in India, a real moat. But production was guided to the second half of 2026 at FID \cite{amgreenfid2024}; shipments now start in 2028. Even the best-capitalised project slipped.\\
\textbf{ReNew} & Listed in New York; Canada's CPPIB holds 31.3\% and ADIA 23.8\% \cite{renew2025} & No recent hydrogen milestone that I could verify; I treat its hydrogen plans as dormant \confL & 10.3 GW installed. A take-private bid at \usd{}8.15 a share collapsed when Masdar withdrew (Dec 2025); shares fell about 28\% to \usd{}5.26 \cite{renew2025} & Size does not guarantee a hydrogen business. Capital structure and investor patience matter as much as assets.\\
\textbf{Avaada} & Founded by Vineet Mittal; raised \usd{}1.07 bn in April 2023, led by Brookfield's transition fund, explicitly for green hydrogen and ammonia \cite{avaada2023} & Bid in five SECI green-ammonia auctions and came second or third in all of them, often by ₹0.25/kg \cite{seci2025} & About 4 GW of renewables at the time of the raise, targeting 11 GW by 2026 \cite{avaada2023} \prj & Raising money early is not the same as winning. Avaada lost every auction to a slightly lower bid, which is either discipline or a cost gap. Either way, capital alone is not an edge.\\
\bottomrule
\end{xltabular}}

\subsection{The technology companies: India and the world}

{\tablefont
\begin{xltabular}{\textwidth}{@{}>{\RaggedRight\arraybackslash}p{18mm} L L L L@{}}
\caption{Technology start-ups and manufacturers. All figures \rep\ unless tagged otherwise.}\\
\toprule
\thead{Company} & \thead{Founding and capital} & \thead{Product and inflection} & \thead{Where it stands (2025–26)} & \thead{Mistake or lesson}\\ \midrule
\endfirsthead
\toprule
\thead{Company} & \thead{Founding and capital} & \thead{Product and inflection} & \thead{Where it stands (2025–26)} & \thead{Mistake or lesson}\\ \midrule
\endhead
\textbf{Ohmium} (US HQ, India factory) & \usd{}45 m Series B led by Fenice \cite{ohmiumb}; \usd{}250 m Series C led by TPG Rise Climate, April 2023 \cite{ohmiumc2023} & PEM electrolysers built in Bengaluru, scaling from 500 MW to 2 GW a year \cite{ohmiumb} \prj; won 137 MW of India's manufacturing incentive \cite{secipli2024} & I could not verify recent orders or financial results \confL & Raised heavily into a market that then shrank. Watch whether its factory runs near capacity.\\
\textbf{Newtrace} (Bengaluru) & Founded 2021 by Prasanta Sarkar and Rochan Sinha; \usd{}6.3 m pre-Series A in March 2026, led by HDFC Bank and Mitsui Sumitomo Insurance VC \cite{newtrace2026} & Sells electrodes (``Voltagen'') into other makers' alkaline stacks rather than building whole systems \cite{newtrace2026} & Pilot manufacturing; no named customers yet \cite{newtrace2026} & A deliberately capital-light ``picks and shovels'' model: it wins if \emph{any} alkaline maker wins.\\
\textbf{Electric Hydrogen} (US) & \usd{}380 m Series C (Oct 2023), the first hydrogen-technology company valued at \usd{}1 bn \cite{eh2gv2023} & Sells whole 100 MW plants (about 50 t of hydrogen a day), not stacks; Fortescue framework for 1 GW; claims an order book above 5 GW \cite{eh2gv2023} & In 2025 partnered with Titan for modular, factory-built plants \cite{eh2titan2025} & Attacked the real cost problem, the balance of plant, rather than the stack alone. Order-book claims still need converting into delivered plants.\\
\textbf{Hysata} (Australia) & \usd{}111.3 m Series B (May 2024), co-led by bp Ventures and Templewater, plus POSCO and others \cite{hysata2024} & Capillary-fed electrolyser claiming efficiency ``well above'' IRENA's 2050 target; no number published in the release \cite{hysata2024} & Scaling manufacturing towards gigawatt scale \cite{hysata2024} \prj & A genuine efficiency bet. Until independent kWh/kg data is public, treat the claim as unproven \confL.\\
\textbf{Enapter} (Germany, AEM) & Listed company & Pioneer of small, modular AEM electrolysers & 2025 revenue €22.1 m, EBITDA $-$€18.1 m, including €4.5 m for \emph{returned stacks} \cite{enapter2026} & The clearest evidence that AEM durability is not yet solved. Directly relevant to Dossier 1's AEM option.\\
\bottomrule
\end{xltabular}}

\subsection{The cautionary tale: Plug Power}

Plug Power is the case every hydrogen founder should study, because it did almost everything the market rewards in a boom and paid for all of it in the bust.

\begin{figure}[H]
\centering
\begin{tikzpicture}[x=1mm,y=1mm,
  ev/.style={anchor=north,text width=31mm,align=center,font=\scriptsize}]
\draw[rule,line width=2pt] (0,0) -- (160,0);
\foreach \x/\yr/\c/\body in {
 10/{Nov 2023}/audit/{Going-concern warning; cash down from about \$2 bn to \$0.5 bn; shares fell about 29\% before the open \cite{plugeagle2023}},
 45/2024/audit/{Revenue \$629 m; net loss above \$2.1 bn \cite{plugfool2026}},
 80/2025/brass/{``Project Quantum Leap'' to cut over \$200 m of costs; shares outstanding up 135\% in three years \cite{plugfool2026}},
 114/{Q4 2025}/sgreen/{First positive gross margin in years (2.4\%); 2025 revenue about \$710 m; operating cash burn \$536 m \cite{plug2026}},
 150/2028/sviolet/{Target: full profitability \cite{plug2026}}}{
 \node[circle,fill=\c,minimum size=3.6mm,inner sep=0] at (\x,0) {};
 \node[font=\display\fontsize{10}{12}\selectfont,text=\c,anchor=south] at (\x,3) {\yr};
 \node[ev] at (\x,-4) {\body};
}
\end{tikzpicture}
\caption{Plug Power, 2023–2028. All figures \rep; the 2028 target is \prj.}
\end{figure}

What went wrong, in plain words. Plug sold equipment and hydrogen to win volume before it could make either at a profit; the company itself blamed ``unprecedented hydrogen supply challenges'' when its 2023 sales disappointed \cite{plugeagle2023}. It funded the gap by selling new shares, so the share count rose 135\% in three years and the stock lost about 85\% \cite{plugfool2026}. It is now cutting costs and selling assets to reach profit by 2028 \cite{plug2026}. \confH

\begin{finding}[Five rules I take from the ten cases]
\begin{enumerate}[leftmargin=1.4em,label={\color{audit}\arabic*}]
\item \textbf{Cash flow first, hydrogen second.} ACME and Greenko built profitable power businesses before betting on molecules. The pure hydrogen plays (Plug, Enapter) burned cash.
\item \textbf{Never sell below cost to buy growth.} Plug's volume came with negative margins and a near-death experience.
\item \textbf{Contracts beat capital.} Avaada raised \usd{}1 bn and still lost every auction. AM Green's turning point was a signed offtake, not a funding round.
\item \textbf{Durability is the hidden cost of new technology.} Enapter's returned stacks are a real charge on real accounts.
\item \textbf{Even the best slip.} AM Green's first exports moved from 2026 to 2028. Plan for delay as the base case, not the risk case.
\end{enumerate}
\confM
\end{finding}
